\chapter{Preference Optimization Variants}
\label{ch:po-variants}

This chapter covers the family of methods that extend or replace DPO with different objectives, data assumptions, or architectural trade-offs. Each addresses a specific limitation of standard offline DPO: distribution shift (Online DPO), the need for paired data (KTO), overfitting to noisy labels (IPO), reference model memory cost (ORPO), or training complexity (Best-of-N).

\section{Online DPO}
\label{sec:online-dpo}

\subsection{Motivation}
\label{motivation}

Standard DPO’s primary limitation: the preference data was generated by a \emph{different} model (often an older checkpoint or even a different model family). As training progresses, the policy generates text that looks nothing like the training pairs $\rightarrow$ the loss is optimizing on an irrelevant distribution.

\textbf{Online DPO solution}~\cite{guo2024direct}: Generate fresh preference pairs from the \emph{current} policy at every step, judge them with a reward model, then apply the DPO loss.

\subsection{Algorithm}
\label{algorithm}

\begin{enumerate}
  \item Generate $K$ responses per prompt from current $\pi_\theta$
  \item Score all responses with reward model $r_\phi$
  \item Create pairs: highest-scoring = chosen, lowest-scoring = rejected
  \item Apply DPO loss on these fresh pairs
  \item Repeat (new generation every step)
\end{enumerate}

\begin{intuitionbox}[Online DPO = Best of Both Worlds]
\begin{itemize}
  \item From DPO: simple supervised loss, no value function, no GAE, stable optimization
  \item From PPO: on-policy data, self-improvement beyond dataset, no distribution shift
  \item Key difference from GRPO: uses DPO loss (pair-based) instead of PPO loss (per-sample advantage)
\end{itemize}

\textbf{Trade-off}: Needs a reward model (DPO doesn’t), but no value head (PPO does). Middle ground complexity.
\end{intuitionbox}


\subsection{TRL Implementation}
\label{trl-implementation}

The following shows a minimal working example using HuggingFace TRL.

\begin{lstlisting}[style=pythonstyle]
from trl import OnlineDPOConfig, OnlineDPOTrainer
from transformers import AutoModelForCausalLM, AutoModelForSequenceClassification

model = AutoModelForCausalLM.from_pretrained("meta-llama/Llama-3.1-8B-Instruct",
    torch_dtype=torch.bfloat16)
reward_model = AutoModelForSequenceClassification.from_pretrained(
    "RLHFlow/ArmoRM-Llama3-8B-v0.1", torch_dtype=torch.bfloat16)

online_dpo_config = OnlineDPOConfig(
    output_dir="./online_dpo_output",
    learning_rate=5e-7,
    beta=0.1,                    # DPO beta (same meaning as standard DPO)
    num_generations=4,           # K responses per prompt
    per_device_train_batch_size=4,
    gradient_accumulation_steps=4,
    max_new_tokens=512,
    temperature=0.7,
    bf16=True,
    num_train_epochs=1,
    logging_steps=10,
)

trainer = OnlineDPOTrainer(
    model=model,
    reward_model=reward_model,
    args=online_dpo_config,
    train_dataset=prompt_dataset,
    tokenizer=tokenizer,
)
trainer.train()
\end{lstlisting}

\subsection{Online DPO vs Offline DPO vs PPO}
\label{online-dpo-vs-offline-dpo-vs-ppo}

\begin{tabular}{@{}llp{4.5cm}lp{4.5cm}@{}}
\toprule
 & \textbf{Data} & \textbf{Models} & \textbf{Loss} & \textbf{Best For} \\
\midrule
Offline DPO & Static pairs & 2 (policy + reference) & DPO & Quick alignment, limited compute \\
Online DPO & Fresh from $\pi_\theta$ & 3 (policy + reference + reward model) & DPO & When DPO plateaus, need exploration \\
PPO & Fresh from $\pi_\theta$ & 4 (policy + reference + reward model + value head) & PPO clip & Max quality, complex reasoning \\
\bottomrule
\end{tabular}

\section{KTO --- Kahneman-Tversky Optimization}
\label{sec:kto}

\subsection{Motivation}
\label{motivation-1}

DPO requires \emph{paired} preferences: for the same prompt, you need both a good and bad response. In practice, most feedback is \emph{unpaired}: users give thumbs up/down on individual responses, with no matched pair.

\textbf{KTO’s insight}~\cite{ethayarajh2024kto}: Use prospect theory (from behavioral economics). Humans feel losses more strongly than gains. A ``thumbs down'' should produce a stronger gradient than a ``thumbs up.''

\subsection{Loss Function}
\label{loss-function}

\begin{equation}
\boxed{\mathcal{L}_\text{KTO} = \mathbb{E}_{y_w}\left[\lambda_w (1 - v(x, y_w))\right] + \mathbb{E}_{y_l}\left[\lambda_l \cdot v(x, y_l)\right]}
\end{equation}
 where $v(x,y) = \sigma\left(\beta \log\frac{\pi_\theta(y|x)}{\pi_\text{ref}(y|x)} - z_\text{ref}\right)$, and $z_\text{ref}$ is the expected KL divergence (a running baseline).

\begin{intuitionbox}[KTO Intuition via Prospect Theory]
\textbf{Desirable responses} ($y_w$): The model gets ``utility'' from increasing their probability. But with diminishing returns --- once it’s already quite likely, don’t push harder.

\textbf{Undesirable responses} ($y_l$): Loss aversion means the penalty for generating bad text is weighted more strongly than the reward for good text. Default: $\lambda_l = 1.0$, $\lambda_w = 1.0$, but you can set $\lambda_l > \lambda_w$.

\textbf{Key advantage}: Each training example is independent! No need to find matched pairs. Can use thumbs-up/down data directly.
\end{intuitionbox}

\begin{examplebox}[KTO Data Format]
Unlike DPO which needs: \verb|{"prompt": ..., "chosen": ..., "rejected": ...}|

KTO only needs: \verb|{"prompt": ..., "completion": ..., "label": true/false}|

This means you can use:

\begin{itemize}
  \item Thumbs up/down from production traffic
  \item Upvotes/downvotes from forums
  \item Human ratings binarized (4--5 stars = good, 1--2 = bad)
  \item Any per-response quality signal
\end{itemize}
\end{examplebox}

\subsection{TRL Implementation}
\label{trl-implementation-1}

The following shows a minimal working example using HuggingFace TRL.

\begin{lstlisting}[style=pythonstyle]
from trl import KTOConfig, KTOTrainer

# Dataset format: {"prompt": str, "completion": str, "label": bool}
# label=True for desirable, label=False for undesirable
kto_dataset = [
    {"prompt": "What's 2+2?", "completion": "The answer is 4.", "label": True},
    {"prompt": "What's 2+2?", "completion": "It might be 5.", "label": False},
]

kto_config = KTOConfig(
    output_dir="./kto_output",
    beta=0.1,
    desirable_weight=1.0,        # Weight for good examples
    undesirable_weight=1.0,      # Weight for bad examples (increase for loss aversion)
    learning_rate=5e-7,
    max_length=2048,
    per_device_train_batch_size=4,
    gradient_accumulation_steps=4,
    num_train_epochs=1,
    bf16=True,
)

trainer = KTOTrainer(
    model=model,
    ref_model=ref_model,  # Or None with LoRA
    args=kto_config,
    train_dataset=kto_dataset,
    tokenizer=tokenizer,
)
trainer.train()
\end{lstlisting}

\subsection{When to Choose KTO}
\label{when-to-choose-kto}

\begin{itemize}
  \item You have binary feedback but \emph{not} matched pairs
  \item Production thumbs-up/down data at scale
  \item One class dominates (e.g., 90\% good, 10\% bad) --- KTO handles imbalance better
  \item Rapid iteration with noisy labels (more robust than DPO to noise)
\end{itemize}

\section{IPO --- Identity Preference Optimization}
\label{ipo-identity-preference-optimization}

\subsection{Motivation}
\label{motivation-2}

DPO has a degenerate solution: it can achieve zero loss by making the margin between chosen and rejected \emph{infinitely large}. In practice, this means DPO overfits --- pushing chosen probability to 1 and rejected to 0, memorizing training data.

\textbf{IPO’s fix}~\cite{azar2024general}: Instead of log-sigmoid (which saturates), use a squared loss that targets a \emph{specific} margin. The loss is minimized at a finite gap, not at infinity.

\subsection{Loss Function}
\label{loss-function-1}

\begin{equation}
\boxed{\mathcal{L}_\text{IPO} = \mathbb{E}\left[\left(\log\frac{\pi_\theta(y_w|x)}{\pi_\text{ref}(y_w|x)} - \log\frac{\pi_\theta(y_l|x)}{\pi_\text{ref}(y_l|x)} - \frac{1}{2\beta}\right)^2\right]}
\end{equation}

\begin{intuitionbox}[IPO vs DPO: Regularization Through Target Margin]
DPO: $\sigma(\text{margin}) \to 1$ optimally. Margin $\to \infty$. No natural stopping point.

IPO: Margin $\to \frac{1}{2\beta}$ optimally. Squared loss penalizes \textbf{both} too-small and too-large margins.

Result: IPO is more robust to noisy labels (a mislabeled pair gets bounded influence), and generalizes better because it doesn’t memorize.
\end{intuitionbox}

\subsection{TRL Implementation}
\label{trl-implementation-2}

The following shows a minimal working example using HuggingFace TRL.

\begin{lstlisting}[style=pythonstyle]
from trl import DPOConfig, DPOTrainer

# IPO is implemented as a DPO loss_type variant in TRL
ipo_config = DPOConfig(
    output_dir="./ipo_output",
    beta=0.1,
    loss_type="ipo",             # The key difference!
    learning_rate=5e-7,
    max_length=2048,
    per_device_train_batch_size=4,
    gradient_accumulation_steps=8,
    bf16=True,
    num_train_epochs=1,
)

trainer = DPOTrainer(
    model=model, ref_model=None, args=ipo_config,
    train_dataset=pref_dataset, tokenizer=tokenizer, peft_config=lora_config,
)
trainer.train()
\end{lstlisting}

\subsection{When to Choose IPO over DPO}
\label{when-to-choose-ipo-over-dpo}

\begin{itemize}
  \item Noisy preference data (crowdsourced, AI-judged with errors)
  \item Observing DPO overfitting (train loss $\to$ 0 but eval degrades)
  \item Want more conservative, robust alignment
  \item Multiple epochs needed (DPO degrades after epoch 1; IPO is more stable)
\end{itemize}

\section{ORPO --- Odds Ratio Preference Optimization}
\label{sec:orpo}

\subsection{Motivation}
\label{motivation-3}

All methods so far need a reference model --- either as a separate copy (doubles memory) or implicitly via LoRA. ORPO~\cite{hong2024orpo} eliminates the reference entirely by combining SFT and preference alignment in a single loss.

\textbf{Key insight}: Use the \emph{odds ratio} of generating chosen vs rejected as the preference signal. The SFT component naturally prevents collapse (no need for KL regularization).

\subsection{Loss Function}
\label{loss-function-2}

\begin{equation}
\boxed{\mathcal{L}_\text{ORPO} = \underbrace{\mathcal{L}_\text{SFT}(y_w)}_{\text{standard NLL on chosen}} - \lambda \cdot \underbrace{\log\sigma\left(\log\frac{\text{odds}_\theta(y_w|x)}{\text{odds}_\theta(y_l|x)}\right)}_{\text{preference alignment via odds ratio}}}
\end{equation}
 where $\text{odds}_\theta(y|x) = \frac{P_\theta(y|x)}{1 - P_\theta(y|x)}$.

\begin{intuitionbox}[ORPO: SFT + Alignment in One Shot]
\textbf{SFT term}: Trains the model to generate the chosen response well (standard language modeling).

\textbf{Odds ratio term}: Additionally pushes the model to prefer chosen over rejected. The odds ratio is a natural contrast that doesn’t require a reference model.

\textbf{Why no reference needed?}: The SFT loss already anchors the model to reasonable text. It serves the same role as KL-to-reference in other methods. One model, one forward pass, one loss. 50\% less memory!
\end{intuitionbox}

\subsection{TRL Implementation}
\label{trl-implementation-3}

The following shows a minimal working example using HuggingFace TRL.

\begin{lstlisting}[style=pythonstyle]
from trl import ORPOConfig, ORPOTrainer

orpo_config = ORPOConfig(
    output_dir="./orpo_output",
    beta=0.1,                    # Odds ratio weight (lambda)
    learning_rate=5e-7,
    max_length=2048,
    per_device_train_batch_size=2,
    gradient_accumulation_steps=8,
    bf16=True,
    num_train_epochs=1,
    gradient_checkpointing=True,
)

trainer = ORPOTrainer(
    model=model,                 # No ref_model needed!
    args=orpo_config,
    train_dataset=pref_dataset,  # Same format as DPO: prompt/chosen/rejected
    tokenizer=tokenizer,
    peft_config=lora_config,
)
trainer.train()
\end{lstlisting}

\subsection{When to Choose ORPO}
\label{when-to-choose-orpo}

\begin{itemize}
  \item Memory-constrained: can’t afford reference model copy (saves 70--140GB for 70B)
  \item Starting from base model (not SFT’d yet) --- ORPO does SFT simultaneously
  \item Want simplest possible pipeline: one model, one loss, one training run
  \item Good preference data available from the start
\end{itemize}

\begin{warningbox}[ORPO Limitations]
\begin{itemize}
  \item Less studied than DPO/PPO --- fewer proven recipes at 70B+ scale
  \item The SFT component means it needs high-quality chosen responses (not just relative preference)
  \item Harder to debug: two loss components can conflict
\end{itemize}
\end{warningbox}

\begin{keybox}[See Also: SimPO]
\textbf{SimPO}~\cite{meng2024simpo} is another reference-free preference method that uses length-normalized log-probability as an implicit reward, eliminating the reference model entirely. It is covered in Section~\ref{sec:simpo} alongside other DPO extensions due to its shared reference-free philosophy.
\end{keybox}

\section{Best-of-N Sampling (Rejection Sampling)}
\label{best-of-n-sampling-rejection-sampling}

\subsection{Motivation}
\label{motivation-4}

Sometimes the simplest approach wins. Best-of-N~\cite{nakano2021webgpt} requires \emph{no training at all} during the RL phase --- just generate multiple candidates and pick the best one.

\subsection{Algorithm}
\label{algorithm-1}

\begin{enumerate}
  \item For each prompt, generate $N$ responses from the policy (typically $N = 4$--$64$)
  \item Score all responses with a reward model
  \item Select the highest-scoring response
  \item (Optional) Use selected responses as SFT data for the next iteration
\end{enumerate}

\begin{equation}
\boxed{\text{Best-of-N response}: \quad y^* = \arg\max_{y_i \sim \pi_\theta(\cdot|x)} r_\phi(x, y_i)}
\end{equation}

\begin{intuitionbox}[Why Best-of-N is a Legitimate ``RL'' Method]
\textbf{At inference time}: Best-of-N improves output quality without changing model weights. With $N=64$, win-rate improves 10--20\% over greedy --- sometimes matching or exceeding PPO.

\textbf{As a training method} (Rejection Sampling Fine-Tuning / RFT):

\begin{enumerate}
  \item Generate many responses, select best ones
  \item SFT on the selected responses
  \item Repeat (iterative refinement)
\end{enumerate}

This is how many production models are trained: simpler than PPO, almost as effective, completely stable.

\textbf{Theoretical connection}~\cite{gao2023scaling}: Best-of-N implements an implicit KL-constrained policy: $\pi_\text{BoN}(y|x) \propto \pi_\theta(y|x)^{1-1/N} \cdot r(x,y)^{1/N}$.
\end{intuitionbox}

\subsection{TRL Implementation}
\label{trl-implementation-4}

The following shows a minimal working example using HuggingFace TRL.

\begin{lstlisting}[style=pythonstyle]
from transformers import pipeline
import numpy as np

# Inference-time Best-of-N (manual implementation)
gen_pipeline = pipeline("text-generation", model=model, tokenizer=tokenizer)

def best_of_n(prompt, n=16, temperature=0.8):
    """Generate N candidates and return the highest-reward one."""
    candidates = gen_pipeline(
        prompt, num_return_sequences=n,
        temperature=temperature, do_sample=True, max_new_tokens=512,
    )
    scores = [reward_model.score(prompt, c["generated_text"]) for c in candidates]
    return candidates[np.argmax(scores)]["generated_text"]

best_response = best_of_n(prompt, n=16)

# Training: Rejection Sampling Fine-Tuning (RFT)
from trl import SFTConfig, SFTTrainer

# Step 1: Generate and filter
all_responses = []
for prompt in prompts:
    candidates = [generate(prompt, temp=0.9) for _ in range(16)]
    scores = [reward_model.score(prompt, c) for c in candidates]
    best_idx = np.argmax(scores)
    if scores[best_idx] > threshold:  # Quality gate
        all_responses.append({"prompt": prompt, "completion": candidates[best_idx]})

# Step 2: SFT on best responses
sft_config = SFTConfig(output_dir="./rft_output", learning_rate=2e-5, num_train_epochs=2, max_seq_length=2048)
trainer = SFTTrainer(model=model, args=sft_config, train_dataset=all_responses, tokenizer=tokenizer)
trainer.train()
# Step 3: Repeat from Step 1 with updated model (iterative RFT)
\end{lstlisting}

\subsection{Scaling Laws for Best-of-N}
\label{scaling-laws-for-best-of-n}

\begin{tabular}{@{}lp{3.5cm}p{3.5cm}p{6cm}@{}}
\toprule
$N$ & \textbf{Quality Gain} & \textbf{Cost} & \textbf{Notes} \\
\midrule
1 & Baseline & $1\times$ & Standard sampling \\
4 & +5--8\% win-rate & $4\times$ & Minimum useful. Good cost/quality ratio \\
16 & +10--15\% win-rate & $16\times$ & Strong. Often matches PPO quality \\
64 & +15--20\% win-rate & $64\times$ & Diminishing returns start \\
256 & +18--22\% win-rate & $256\times$ & Only for critical applications \\
\bottomrule
\end{tabular}

\begin{keybox}[Best-of-N as Baseline]
Always compare your RL method against Best-of-N with the same compute budget. If PPO with 64 GPU-hours doesn’t beat Best-of-N with 64 GPU-hours of generation, your PPO has a bug.
\end{keybox}

\section{Summary: Choosing an Alignment Method}
\label{sec:method-comparison}

We have now surveyed the full landscape of preference optimization and RL-based alignment methods. This section consolidates the key trade-offs into a single reference to help practitioners select the right approach for their constraints.

\begin{table}[ht!]
\centering
\caption{Cross-method comparison of alignment approaches.}
{\footnotesize
\begin{tabular}{@{}llllll@{}}
\toprule
\textbf{Method} & \textbf{Models} & \textbf{Data} & \textbf{Compute} & \textbf{Stability} & \textbf{Best For} \\
\midrule
PPO & 4 & Online (gen) & Very high & Low & Max quality, complex reasoning \\
GRPO & 2 (no critic) & Online (gen) & High & Medium & Math/code (verifiable rewards) \\
DPO & 2 & Offline pairs & Low & High & Style/safety, limited compute \\
Online DPO & 3 & Online (gen) & Medium & Medium-High & DPO without distribution shift \\
KTO & 2 & Unpaired binary & Low & High & Production feedback, thumbs up/down \\
IPO & 2 & Offline pairs & Low & Very high & Noisy labels, anti-overfitting \\
ORPO & 1 & Offline pairs & Very low & High & Memory-limited, SFT+align combined \\
Best-of-N & 1+RM & Online (gen) & Medium & Perfect & Strong baseline, data generation \\
\bottomrule
\end{tabular}
}
\end{table}

\begin{figure}[ht!]
\centering
\includegraphics[width=0.85\textwidth]{figures/fig_029_fig29.png}
\caption{Approximate quality vs.~compute frontier. Methods above the SFT ceiling line improve beyond what supervised fine-tuning alone achieves. Position is illustrative and model-dependent.}
\end{figure}

\begin{keybox}[Decision Tree: Which Method to Use?]
\begin{enumerate}
  \item \textbf{Do you have verifiable rewards?} (math/code) $\rightarrow$ \textbf{GRPO}
  \item \textbf{Do you need max quality on complex tasks?} $\rightarrow$ \textbf{PPO}
  \item \textbf{Do you have paired preferences?} $\rightarrow$ \textbf{DPO} (or IPO if noisy)
  \item \textbf{Only unpaired binary feedback?} $\rightarrow$ \textbf{KTO}
  \item \textbf{Memory-limited, starting from base model?} $\rightarrow$ \textbf{ORPO}
  \item \textbf{DPO plateauing, want on-policy?} $\rightarrow$ \textbf{Online DPO}
  \item \textbf{Need a strong baseline quickly?} $\rightarrow$ \textbf{Best-of-N / RFT}
\end{enumerate}
\end{keybox}

